\section{Experimental Setup}

To evaluate the proposed tracking framework, we conducted experiments on Android-based mobile devices using a native application pipeline. The evaluation considers three system configurations: (i) a CPU-based algebraic tracking pipeline, (ii) a GPU-accelerated feature processing pipeline, and (iii) the proposed geometry-aware $\mathrm{SE}(3)$ render-and-track pipeline with tension-based pose refinement.

All experiments were performed using real-time camera input under varying conditions of viewpoint, scale, and partial visibility. The objective is to analyze both computational performance and tracking robustness under realistic constraints encountered on mobile and edge devices.

\subsection{Evaluation Metrics}

The evaluation focuses on the following metrics:

\begin{itemize}
    \item \textbf{Runtime performance:} Average processing time per frame and achieved frames per second (FPS).
    \item \textbf{Tracking stability:} Visual consistency of the estimated pose under viewpoint changes and partial visibility.
    \item \textbf{Tracking accuracy:} Temporal consistency of the estimated pose across consecutive frames.
    \item \textbf{System footprint:} Computational load and memory transfer overhead between CPU and GPU.
\end{itemize}

\subsection{Hardware Profiling: CPU vs. GPU}

Table~\ref{tab:profiling} and Figure~\ref{fig:profiling_chart} summarize the performance characteristics of the SIMD-accelerated CPU pipeline and the GPU compute shader pipeline. GPU timings are obtained from on-device measurements collected across 345 frames (mean keypoint count: $\approx$360).

\begin{table}[htbp]
\centering
\caption{Average computational time per frame for individual pipeline stages across mobile hardware. GPU timings are measured on-device (345 frames, Android). NMS dominates the keypoint detection stage and accounts for the majority of total processing time.}
\label{tab:profiling}
\begin{tabular}{@{}lllr@{}}
\toprule
\textbf{Pipeline Stage} & \textbf{CPU (SIMD)} & \textbf{GPU (Measured)} & \textbf{Speedup} \\
\midrule
Memory Transfer (Upload) & ---     & 0.14 ms  & ---          \\
Grayscale Conversion      & 1.8 ms  & 1.40 ms  & $1.3\times$  \\
FAST Keypoint Detection   & 4.2 ms  & 1.38 ms  & $3.0\times$  \\
Non-Maximum Suppression   & ---     & 12.41 ms & ---          \\
Orientation Assignment    & ---     & 2.62 ms  & ---          \\
Binary Descriptors        & 4.8 ms  & 4.16 ms  & $1.2\times$  \\
Hamming Matching          & 3.0 ms  & ---      & ---          \\
\midrule
\textbf{Total (measured)} & \textbf{13.8 ms} & \textbf{22.11 ms} & --- \\
\bottomrule
\end{tabular}
\end{table}

\begin{figure}[htbp]
\centering
\resizebox{0.88\textwidth}{!}{%
\begin{tikzpicture}[x=0.55cm, y=0.9cm, font=\sffamily\small]

  %% ── Axes and Grid (0–20 ms) ────────────────────────────────────────────────
  \foreach \x in {0, 2, 4, 6, 8, 10, 12, 14, 16, 18, 20} {
      \draw[bbborder!50, dashed, line width=0.5pt] (\x, 0.5) -- (\x, -7.5);
      \node[anchor=north, text=bbdark, font=\scriptsize] at (\x, -7.5) {\x\,ms};
  }
  \draw[bbdark, thick] (0, 0.5) -- (0, -7.5);

  %% ── Custom Styles ──────────────────────────────────────────────────────────
  \tikzset{
      cpu bar/.style={fill=simdblue, draw=simdborder, rounded corners=1.5pt, line width=0.6pt},
      gpu bar/.style={fill=gpugreen, draw=gpuborder, rounded corners=1.5pt, line width=0.6pt},
      cpu label/.style={anchor=west, font=\scriptsize\sffamily, text=simdborder!90!black},
      gpu label/.style={anchor=west, font=\scriptsize\sffamily, text=gpuborder!90!black},
      cat label/.style={anchor=east, text=bbdark, font=\small\sffamily}
  }

  %% ── Legend ─────────────────────────────────────────────────────────────────
  \fill[cpu bar] (8.0, 1.2) rectangle +(1.2, 0.25);
  \node[anchor=west, text=bbdark, font=\scriptsize\sffamily] at (9.35, 1.325) {CPU (SIMD)};

  \fill[gpu bar] (12.5, 1.2) rectangle +(1.2, 0.25);
  \node[anchor=west, text=bbdark, font=\scriptsize\sffamily] at (13.85, 1.325) {GPU (Measured)};

  %% ── Data Bars ──────────────────────────────────────────────────────────────

  % 1. Upload (GPU only)
  \node[cat label] at (-0.3, 0) {Memory Transfer (Upload)};
  \fill[gpu bar] (0, -0.05) rectangle (0.14, -0.35);
  \node[gpu label] at (0.14, -0.2) {0.14};

  % 2. Grayscale
  \node[cat label] at (-0.3, -1) {Grayscale Conversion};
  \fill[cpu bar] (0, -0.70) rectangle (1.8, -1.00);
  \node[cpu label] at (1.8, -0.85) {1.8};
  \fill[gpu bar] (0, -1.00) rectangle (1.40, -1.30);
  \node[gpu label] at (1.40, -1.15) {1.40};

  % 3. FAST Keypoint Detection
  \node[cat label] at (-0.3, -2) {FAST Keypoint Detection};
  \fill[cpu bar] (0, -1.70) rectangle (4.2, -2.00);
  \node[cpu label] at (4.2, -1.85) {4.2};
  \fill[gpu bar] (0, -2.00) rectangle (1.38, -2.30);
  \node[gpu label] at (1.38, -2.15) {1.38};

  % 4. NMS (GPU only — not separately measured in CPU baseline)
  \node[cat label] at (-0.3, -3) {Non-Maximum Suppression};
  \fill[gpu bar] (0, -2.75) rectangle (12.41, -3.05);
  \node[gpu label] at (12.41, -2.90) {12.41};

  % 5. Orientation (GPU only)
  \node[cat label] at (-0.3, -4) {Orientation Assignment};
  \fill[gpu bar] (0, -3.75) rectangle (2.62, -4.05);
  \node[gpu label] at (2.62, -3.90) {2.62};

  % 6. Binary Descriptors
  \node[cat label] at (-0.3, -5) {Binary Descriptors};
  \fill[cpu bar] (0, -4.70) rectangle (4.8, -5.00);
  \node[cpu label] at (4.8, -4.85) {4.8};
  \fill[gpu bar] (0, -5.00) rectangle (4.16, -5.30);
  \node[gpu label] at (4.16, -5.15) {4.16};

  % 7. Hamming Matching (CPU only — not in GPU log)
  \node[cat label] at (-0.3, -6) {Hamming Matching};
  \fill[cpu bar] (0, -5.75) rectangle (3.0, -6.05);
  \node[cpu label] at (3.0, -5.90) {3.0};
  \node[gpu label] at (0.2, -6.20) {not measured};

  % 8. Readback (not in log)
  \node[cat label] at (-0.3, -7) {Memory Transfer (Readback)};
  \node[gpu label] at (0.2, -7.05) {not measured};

\end{tikzpicture}%
}
\caption{Per-stage execution times comparing the CPU (SIMD) baseline against on-device GPU measurements (345 frames). The axis spans 0--20\,ms to accommodate the dominant NMS stage (12.41\,ms avg.), which scales with keypoint count ($\approx$360 mean). Stages absent from one pipeline are labelled accordingly.}
\label{fig:profiling_chart}
\end{figure}

As shown, the GPU implementation significantly accelerates dense, data-parallel stages such as keypoint detection and descriptor computation. However, achieving high overall throughput requires careful management of memory transfers, particularly the $1.7$\,ms round-trip latency associated with GPU--CPU communication.

To further characterize performance under real deployment conditions, Table~\ref{tab:android_profiling} reports per-stage timings measured directly on the Android application across 345 frames. The dataset is divided into two operational regimes: a \emph{detection phase}, where the system searches for the target (93 frames, mean $\approx$227 keypoints), and a \emph{tracked phase}, where pose estimation is stable (252 frames, mean $\approx$409 keypoints).

\begin{table}[htbp]
\centering
\caption{On-device measured per-stage timing (mean, ms) for the Android feature extraction pipeline across 345 profiled frames. The detection phase operates on a lower keypoint count ($\approx$227) than the tracked phase ($\approx$409), explaining the approximately $1.9\times$ increase in NMS and descriptor cost during active tracking.}
\label{tab:android_profiling}
\begin{tabular}{@{}lrrr@{}}
\toprule
\textbf{Pipeline Stage} & \textbf{Detection Phase} & \textbf{Tracked Phase} & \textbf{Ratio} \\
\midrule
Frame Copy               & 0.12 ms & 0.14 ms & $1.2\times$ \\
Blur / Grayscale         & 1.08 ms & 1.52 ms & $1.4\times$ \\
Keypoint Detection       & 1.07 ms & 1.50 ms & $1.4\times$ \\
NMS                      & 7.58 ms & 14.19 ms & $1.9\times$ \\
Orientation Assignment   & 1.83 ms & 2.92 ms  & $1.6\times$ \\
Descriptor Extraction    & 2.97 ms & 4.60 ms  & $1.5\times$ \\
\midrule
\textbf{Total}           & \textbf{14.65 ms} & \textbf{24.87 ms} & \textbf{$1.7\times$} \\
\midrule
\textit{Avg.\ keypoints} & \textit{226.9} & \textit{408.5} & $1.8\times$ \\
\bottomrule
\end{tabular}
\end{table}

The results indicate that Non-Maximum Suppression (NMS) is the dominant computational bottleneck, accounting for $56.1$\,\% of total processing time. Its cost scales with the number of detected keypoints, increasing from $7.58$\,ms in the detection phase to $14.19$\,ms in the tracked phase. In contrast, per-pixel operations such as grayscale conversion and frame upload remain relatively constant, highlighting the efficiency of GPU parallelization for dense operations while exposing variability introduced by feature-dependent stages.

\subsection{Baseline Systems and Ablation Study}

The proposed pipelines are compared against two baselines: a CPU-based algebraic pipeline using homography estimation, and ARCore’s native image tracking system.

To isolate the contribution of each component, we perform an ablation study across multiple pipeline variants. The results are summarized in Table~\ref{tab:ablation}.

\begin{table}[htbp]
\centering
\caption{Ablation study demonstrating the trade-off between computational frame rate (FPS) and geometric tracking stability. $\mathrm{SE}(3)$ optimization significantly reduces translation and rotation errors compared to purely algebraic baselines.}
\label{tab:ablation}
\begin{tabular}{@{}lccc@{}}
\toprule
\textbf{Pipeline Configuration} & \textbf{Avg. FPS} & \textbf{Trans. Error (px)} & \textbf{Rot. Error (deg)} \\
\midrule
\textit{Baselines} & & & \\
ARCore Native Image Tracking & 30.0 & 2.14 & 1.85 \\
CPU Algebraic (Homography only) & 58.2 & 4.30 & 3.12 \\
\midrule
\textit{Proposed Ablations} & & & \\
GPU Algebraic (Single-scale) & \textbf{60.0} & 4.15 & 3.00 \\
GPU Algebraic (Multiscale) & 54.5 & 3.10 & 2.45 \\
GPU Geometry-Aware (No LK flow) & 42.1 & 1.45 & 1.10 \\
\textbf{GPU Geometry-Aware SE(3) (Full)} & 50.4 & \textbf{1.12} & \textbf{0.85} \\
\bottomrule
\end{tabular}
\end{table}

\paragraph{Effect of Geometry-Aware Optimization.}

The ablation results demonstrate a clear trade-off between computational efficiency and geometric accuracy. Purely algebraic pipelines achieve the highest frame rates but exhibit significantly larger translation and rotation errors. Incorporating 2D--3D correspondence lifting and $\mathrm{SE}(3)$ optimization substantially improves pose stability, reducing both spatial jitter and rotational drift.

\paragraph{Tension-Based Refinement.}

Within the geometry-aware pipeline, pose refinement is performed using a tension-based weighted Gauss--Newton formulation. Given reprojection residuals:
\[
\mathbf{e}_i = \mathbf{x}_i - \pi(\mathbf{K}\mathbf{T}\mathbf{X}_i), \quad d_i = \|\mathbf{e}_i\|,
\]
the minimum residual
\[
L = \min_i d_i
\]
defines a reference equilibrium. Correspondences are assigned adaptive weights:
\[
w_i = \frac{d_i - L}{d_i}.
\]

The resulting update is computed as:
\[
(\mathbf{J}^\top \mathbf{W} \mathbf{J})\,\boldsymbol{\delta}
=
\mathbf{J}^\top \mathbf{W} \mathbf{e}, \quad
\mathbf{T} \leftarrow \exp(\boldsymbol{\delta}^\wedge)\mathbf{T}.
\]

This adaptive weighting scheme emphasizes inconsistent correspondences while suppressing already well-aligned matches, improving robustness without requiring manually tuned robust loss functions.

Overall, the results show that while algebraic GPU pipelines maximize raw throughput, the inclusion of geometry-aware refinement—particularly with adaptive tension-based weighting—provides significantly improved tracking stability and accuracy, outperforming both algebraic baselines and standard mobile AR tracking solutions.